\chapter{Luồng tương tác giữa các tác tử}
\label{chap:workflow}

\section{Luồng xử lý tuần tự và song song (Sequential \& Parallel Workflow)}
\label{sec:sequential-workflow}

Luồng tương tác giữa các tác tử trong một lượt hội chẩn (\emph{consultation turn}) tuân thủ quy trình 11~bước chặt chẽ, kết hợp kiểm định tuần tự ở đầu vào với suy diễn song song ở tầng chuyên môn:

\begin{enumerate}
  \item \textbf{Bước 1}: Người dùng gửi các chỉ số sức khỏe, tệp bệnh án đính kèm (nếu có) và câu hỏi tự nhiên qua giao diện Web.
  \item \textbf{Bước 2}: Tầng Cổng API (\texttt{routes\_chat.py}) tiếp nhận yêu cầu, tạo mã phiên và chuyển luồng điều khiển cho Tác tử Điều phối (\agent{Orchestrator}).
  \item \textbf{Bước 3}: Tác tử Dữ liệu Bệnh nhân (\agent{DocumentAgent} và \agent{ConversationAgent}) bóc tách chỉ số, chuẩn hóa đơn vị đo lường, kiểm tra tính hợp lệ của \texttt{PatientProfile} và phối hợp với \agent{ClarificationAgent} kiểm tra phạm vi sàng lọc.
  \item \textbf{Bước 4}: Nếu dữ liệu không hợp lệ, thiếu chỉ số bắt buộc (chưa có cả đường huyết/HbA1c lẫn cặp tuổi + BMI) hoặc gặp điều kiện chặn luồng (\texttt{blocking=True} như đang mang thai, bệnh nhi), Tác tử Điều phối dừng sớm luồng suy diễn và trả về thông báo yêu cầu làm rõ / bổ sung dữ liệu.
  \item \textbf{Bước 5}: Tác tử Dự đoán Học sâu (\agent{PredictionAgent}) thực thi song song cặp mô hình đã huấn luyện (\texttt{nhanes\_mlp.pt} + \texttt{nhanes\_gbdt.joblib}) cùng mô hình đối chứng quy mô lớn BRFSS (\texttt{brfss\_mlp.pt} + \texttt{brfss\_gbdt.joblib}).
  \item \textbf{Bước 6}: Tác tử Truy xuất Tri thức (\agent{KnowledgeAgent}) đối chiếu xét nghiệm với ngưỡng ADA 2024 / Bộ Y tế, đồng thời truy xuất đường đi trên Đồ thị Tri thức Neo4j và các đoạn hướng dẫn từ Qdrant/FTS5.
  \item \textbf{Bước 7}: Tác tử Luật và Phản biện (\agent{ConsensusCriticAgent}) đánh giá các luật lâm sàng tường minh; nếu phát hiện mâu thuẫn giữa dự báo Học sâu và xét nghiệm sinh hóa, tác tử này phát thông điệp chất vấn (\texttt{CHALLENGE}) để các tác tử liên quan giải trình.
  \item \textbf{Bước 8}: Tác tử Học tăng cường (\agent{TreatmentRLAgent}) ánh xạ trạng thái bệnh nhân vào bảng Q-table (\texttt{rl\_qtable.npz}) và áp dụng Cổng an toàn lâm sàng để chọn khuyến cáo can thiệp kèm danh mục thuốc tham chiếu.
  \item \textbf{Bước 9}: Tác tử Điều phối (\agent{Orchestrator}) thu thập và hợp nhất toàn bộ đầu ra từ các tác tử chuyên môn.
  \item \textbf{Bước 10}: Tác tử Giải thích (\agent{ResponseAgent}) sinh báo cáo hỗ trợ quyết định lâm sàng có cấu trúc 5~phần, trong khi Tác tử Giám sát (\texttt{TraceCollector}) ghi nhận toàn bộ siêu dữ liệu thực thi (\emph{execution metadata}).
  \item \textbf{Bước 11}: Báo cáo hoàn chỉnh kèm vết suy luận và trích dẫn nguồn được trả về cho người dùng.
\end{enumerate}

\section{Biểu đồ tuần tự (Sequence Diagram)}
\label{sec:sequence-diagram}

\figref{fig:sequence-diagram} trình bày biểu đồ tuần tự (\emph{sequence diagram}) mô tả luồng trao đổi thông điệp giữa người dùng, Tác tử Điều phối và các tác tử chuyên môn trong một chu trình hội chẩn chuẩn.

\begin{figure}[H]
\centering
\begin{adjustbox}{max width=\linewidth}
\begin{tikzpicture}[
  font=\scriptsize,
  actor/.style={draw=blue!75!black, fill=blue!8, rounded corners=2pt, minimum width=1.85cm, minimum height=0.8cm, align=center, font=\scriptsize\bfseries},
  msg/.style={-{Stealth[length=2mm,width=1.4mm]}, line width=0.55pt, draw=black!80},
  ret/.style={-{Stealth[length=2mm,width=1.4mm]}, line width=0.55pt, dashed, draw=teal!75!black},
  life/.style={dashed, draw=black!40, line width=0.45pt}
]

% Actors across the top
\node[actor] (u) at (0,0) {Người dùng\\(User)};
\node[actor] (c) at (2.4,0) {Coordinator\\(\agent{Orchestrator})};
\node[actor] (d) at (4.9,0) {Data Agent\\(\agent{Conversation})};
\node[actor] (dl) at (7.4,0) {DL Agent\\(\agent{Prediction})};
\node[actor] (r) at (9.9,0) {Rule/Critic\\(\agent{Consensus})};
\node[actor] (kg) at (12.4,0) {KG/RAG Agent\\(\agent{Knowledge})};
\node[actor] (e) at (14.9,0) {Explanation\\(\agent{Response})};

% Lifelines
\foreach \n in {u,c,d,dl,r,kg,e} {
  \draw[life] (\n.south) -- ++(0,-8.6cm);
}

% Step 1: User -> Coordinator
\draw[msg] (0,-0.9) -- node[above, font=\tiny]{1. Yêu cầu khám / Câu hỏi} (2.4,-0.9);

% Step 2: Coordinator -> Data Agent
\draw[msg] (2.4,-1.55) -- node[above, font=\tiny]{2. \texttt{PROCESS\_TURN} (Kiểm định)} (4.9,-1.55);
\draw[ret] (4.9,-2.15) -- node[above, font=\tiny]{3. Hồ sơ hợp lệ (\texttt{PatientProfile})} (2.4,-2.15);

% Step 3: Coordinator -> DL Agent
\draw[msg] (2.4,-2.85) -- node[above, font=\tiny]{4. \texttt{PREDICT\_RISK}} (7.4,-2.85);
\draw[ret] (7.4,-3.45) -- node[above, font=\tiny]{5. Dự báo DL/ML ($p_{\text{ens}}, \text{RR}$)} (2.4,-3.45);

% Step 4: Coordinator -> KG/RAG Agent
\draw[msg] (2.4,-4.15) -- node[above, font=\tiny]{6. \texttt{RETRIEVE\_EVIDENCE}} (12.4,-4.15);
\draw[ret] (12.4,-4.75) -- node[above, font=\tiny]{7. Bằng chứng Neo4j KG + Qdrant RAG} (2.4,-4.75);

% Step 5: Coordinator -> Rule/Critic Agent (+ Debate + RL)
\draw[msg] (2.4,-5.45) -- node[above, font=\tiny]{8. \texttt{EVALUATE\_CONSENSUS} (+ RL)} (9.9,-5.45);
\draw[ret] (9.9,-6.10) -- node[above, font=\tiny]{9. Kết quả luật, cờ phản biện \& phác đồ} (2.4,-6.10);

% Step 6: Coordinator -> Explanation Agent
\draw[msg] (2.4,-6.85) -- node[above, font=\tiny]{10. \texttt{COMPOSE\_ASSESSMENT}} (14.9,-6.85);
\draw[ret] (14.9,-7.45) -- node[above, font=\tiny]{11. Báo cáo 5 phần có dẫn chứng} (2.4,-7.45);

% Step 7: Coordinator -> User
\draw[ret] (2.4,-8.15) -- node[above, font=\tiny]{12. Trả báo cáo + Vết kiểm toán (\texttt{Trace})} (0,-8.15);

\end{tikzpicture}
\end{adjustbox}
\caption{Biểu đồ tuần tự tương tác giữa các tác tử trong Hệ thống Đa Tác tử Lai.}
\label{fig:sequence-diagram}
\end{figure}

\section{Thuật toán điều phối của Coordinator (Coordinator Workflow)}
\label{sec:coordinator-workflow}

Tác tử điều phối \agent{Orchestrator} hiện thực hóa thuật toán điều khiển 12~bước chuẩn hóa sau đây trong phương thức \texttt{run\_turn()}:
\begin{enumerate}
  \item Tiếp nhận yêu cầu từ người dùng (kèm văn bản, danh sách tệp đính kèm và trạng thái hồ sơ hiện tại).
  \item Khởi tạo mã định danh lượt xử lý (\texttt{trace\_id}) và đối tượng giám sát \texttt{TraceCollector}.
  \item Gọi Tác tử Dữ liệu Bệnh nhân (\agent{DocumentAgent} nếu có tệp, sau đó là \agent{ConversationAgent}) để bóc tách và hợp nhất \texttt{PatientProfile}.
  \item Kiểm tra trạng thái hợp lệ, ý định hội thoại và điều kiện phạm vi qua \agent{ClarificationAgent}.
  \item Nếu hồ sơ chưa đủ chỉ số tối thiểu hoặc gặp điều kiện chặn (\texttt{blocking\_questions}), trả về ngay thông báo hướng dẫn bổ sung dữ liệu.
  \item Kích hoạt Tác tử Dự đoán Học sâu (\agent{PredictionAgent}) với thông điệp \texttt{PREDICT\_RISK}.
  \item Kích hoạt Tác tử Truy xuất Tri thức (\agent{KnowledgeAgent}) với thông điệp \texttt{RETRIEVE\_EVIDENCE}.
  \item Kích hoạt Tác tử Luật và Phản biện (\agent{ConsensusCriticAgent}) với thông điệp \texttt{EVALUATE\_CONSENSUS} (bao gồm việc thực thi các thông điệp \texttt{CHALLENGE} nếu phát hiện mâu thuẫn).
  \item Kích hoạt Tác tử Học tăng cường (\agent{TreatmentRLAgent}) với thông điệp \texttt{RECOMMEND\_TREATMENT} và thu thập toàn bộ đầu ra.
  \item Gọi Tác tử Giải thích (\agent{ResponseAgent}) với thông điệp \texttt{COMPOSE\_ASSESSMENT} để tổng hợp báo cáo 5~phần.
  \item Ghi nhận toàn bộ vết thực thi (\texttt{trace.transcript()}) và biểu đồ tương tác (\texttt{trace.interaction\_graph()}).
  \item Trả về cấu trúc kết quả hoàn chỉnh cho tầng API.
\end{enumerate}

\coderef{code:orchestrator-core} trích dẫn đoạn mã thực tế của phương thức \texttt{Orchestrator.run\_turn()} trong tệp \texttt{src/agents/orchestrator.py}.

\begin{code}{Thuật toán điều phối trung tâm trong \texttt{src/agents/orchestrator.py}}
\label{code:orchestrator-core}
\begin{minted}{python}
# Bước 5 & 6: Kích hoạt song song DL Prediction Agent và KG/RAG Agent
pred_msg, know_msg = await asyncio.gather(
    self.request("PredictionAgent", "PREDICT_RISK",
                 {"profile": profile.model_dump()}, trace),
    self.request("KnowledgeAgent", "RETRIEVE_EVIDENCE",
                 {"profile": profile.model_dump()}, trace),
)
prediction, knowledge = pred_msg.payload, know_msg.payload

# Bước 7: Rule-Based & Consensus Critic Agent kiểm định và tranh biện đa vòng
consensus = await self._run_consensus_debate(profile, prediction, knowledge, trace)

# Bước 8: RL Agent tra cứu chính sách can thiệp + Safety Gate
treat_msg = await self.request("TreatmentRLAgent", "RECOMMEND_TREATMENT", {
    "profile": profile.model_dump(),
    "consensus": consensus,
    "knowledge": knowledge,
}, trace)

# Bước 10: Explanation Agent tổng hợp báo cáo 5 phần có dẫn chứng
resp = await self.request("ResponseAgent", "COMPOSE_ASSESSMENT", {
    "profile": profile.model_dump(), "prediction": prediction,
    "knowledge": knowledge, "consensus": consensus,
    "treatment": treat_msg.payload, "questions": questions,
}, trace)
\end{minted}
\end{code}
